% Λύση του 13ου προβλήματος της ενότητας «Άλυτα Προβλήματα».
\subsection*{Λύση Προβλήματος 13 --- Μέγιστος φωτισμός στο άκρο}

Η λάμπα, το κέντρο του τραπεζιού και ένα σημείο στο άκρο του σχηματίζουν
ορθογώνιο τρίγωνο με κάθετες πλευρές $h$ και $R$. Επομένως,
\[
  d=\sqrt{h^2+R^2}.
\]

Η γωνία $\theta$ σχηματίζεται από την οπτική ακτίνα και την επιφάνεια
του τραπεζιού. Άρα
\[
  \hm\theta=\frac{h}{d}
  =\frac{h}{\sqrt{h^2+R^2}}.
\]
Από τη σχέση
\[
  I=k\frac{\hm\theta}{d^2}
\]
προκύπτει
\begin{align*}
  I(h)
  &=k\frac{h}{\sqrt{h^2+R^2}}
       \frac{1}{h^2+R^2}\\
  &=\frac{kh}{(h^2+R^2)^{3/2}},
  \qquad h>0.
\end{align*}
Επομένως,
\[
  \boxed{I(h)=\frac{kh}{(h^2+R^2)^{3/2}}}.
\]

Παραγωγίζουμε:
\begin{align*}
  I'(h)
  &=k(h^2+R^2)^{-3/2}
    -3kh^2(h^2+R^2)^{-5/2}\\
  &=\frac{k(R^2-2h^2)}
          {(h^2+R^2)^{5/2}}.
\end{align*}

Για $h>0$ ο παρονομαστής και η σταθερά $k$ είναι θετικοί. Άρα το
πρόσημο της $I'(h)$ καθορίζεται από την παράσταση $R^2-2h^2$.
Έχουμε
\[
  I'(h)=0
  \Longleftrightarrow
  R^2-2h^2=0
  \Longleftrightarrow
  h=\frac{R}{\sqrt2}.
\]

Επιπλέον,
\[
\begin{array}{c|ccc}
  h & \left(0,\frac{R}{\sqrt2}\right)
    & \frac{R}{\sqrt2}
    & \left(\frac{R}{\sqrt2},+\infty\right)\\ \hline
  I'(h) & + & 0 & -\\
  I(h) & \nearrow & \text{μέγιστο} & \searrow
\end{array}
\]

Συνεπώς, ο φωτισμός στο άκρο του τραπεζιού γίνεται μέγιστος όταν
\[
  \boxed{h=\frac{R}{\sqrt2}}.
\]
